\section{Causal Inference in NLP}
\label{sec:rw-causal-inference-nlp}

Causal inference methods have been applied to text analysis across a range of settings, and the literature is usefully organized by the role text plays in the causal graph. In \emph{text-as-treatment} settings, features extracted from text (sentiment, framing, topic) are treated as the intervention whose effect on a downstream outcome is estimated. In \emph{text-as-outcome} settings, the text itself is the response variable, and the question is how some non-textual treatment shifts language. In \emph{text-as-confounder} settings, text is used to adjust for unobserved structure that would otherwise bias an estimated effect. \citet{feder2022causalinferenceinnlp} provide a comprehensive survey across these settings; \citet{veitch2019textembeddingscausalinference} propose using text embeddings to adjust for confounding; \citet{keith-etal-2021-text} develop a framework for text as causal mediator, estimating direct and indirect effects of social group signals through language; and \citet{tierney-volfovsky-2021-sensitivity} apply causal mediation analysis on conversation (direct effect) and its content (indirect effect) on depolarisation.

In parallel, causal methods have been applied to understand the behaviour of language models themselves. \citet{vig2020genderbiasLLMs} use causal mediation analysis to localize and diagnose gender bias in language models, decomposing how gendered surface cues propagate through model components into biased outputs. \citet{stolfo2023mechanisticinterpretationarithmeticreasoning} apply mechanistic interpretation via causal mediation to arithmetic reasoning, and \citet{han2024surfacestructurecausalassessment} use causal probing to assess LLM comprehension beyond surface structure.

Our work in Chapter~\ref{sec:causal-inference} contributes to this literature by applying causal inference in a previously unexplored role: as an \emph{annotation-quality diagnostic} rather than as a model-internals interpretability tool. Instead of opening up a single model and asking which neurons mediate a behaviour, we take entire annotation pipelines (human, fine-tuned encoders, zero-shot LLMs, fine-tuned LLMs) as black-box label generators and ask whether the causal structure linking sentiment, topic, and ideology is preserved across them. Divergences in that causal structure (not divergences in F1) are what reveal shortcut learning. To our knowledge, framing causal mediation as a downstream auditor of annotator trustworthiness, rather than as a probe of a fixed model's internals, is a novel use of these methods in NLP.

\subsection{Double Machine Learning}
\label{sec:dml-related-work}

The specific estimator this thesis uses to operationalize the text-as-treatment analysis is Double Machine Learning (DML), introduced by \citet{chernozhukov2018}. DML is a causal inference framework for estimating treatment effects when the number of potential confounders is large relative to sample size. Standard regression of an outcome on a treatment is biased when confounders predict both, and conditioning on every observed covariate becomes statistically unstable as covariate dimensionality grows. DML resolves this by residualizing both the treatment and the outcome against flexible (typically machine-learned) nuisance models of the confounders, and then regressing the residualized outcome on the residualized treatment via cross-fitting and Neyman-orthogonal score functions. The result is a root-$N$ consistent estimate of the treatment effect that remains valid even when the nuisance models are themselves high-dimensional and only converge at slower rates.

DML is well suited to the setting in Chapter~\ref{sec:causal-inference}, where the treatment is a continuous topic-level sentiment score, the outcome is an article-level ideology label, and the confounder set is the high-dimensional topic-presence indicator over hundreds of AllSides topic tags. Methodological details of how we instantiate DML, including the choice of base learners, the cross-fitting protocol, and the extension to mediation analysis, are deferred to Section~\ref{sec:dml}.
